\begin{abstract}
DiCo-NLI, sıralı ifade çiftlerinin eşdeğerlik, ileri gerektirme, geri gerektirme ya da başka bir ilişki olarak etiketlenmesini ister. Ayrıca bir çiftin iki yönü için verilen tahminlerin tersine çevirme eşlemesi altında uyumlu olup olmadığını puanlar. Bir sınıflandırıcının bu cebiri kendiliğinden öğrenmesini beklemek yerine onu çıkarım sırasında zorunlu kılıyoruz: çok dilli bir kodlayıcı iki yönü de puanlar, \emph{ortak çift yönlü çözümleme} ise çift için bir bütün olarak en iyi etiketi seçer. Böylece iki tahmin tasarım gereği tutarlı olur. İngilizce, İspanyolca, Baskça ve karma dilli track'lerde ortak çözümleme, eşli doğruluğu (HardCons) 72 değerlendirmenin tamamında ek eğitim gerektirmeden artırır ve eğitim sırasında uygulanan bir tutarlılık kaybından daha etkilidir. En iyi sistemimiz, NLI ön eğitimli mDeBERTa'nın ortak çözümlemeli 3 tohumlu topluluğudur ve dev kümelerinde ortalama 78.9 ağırlıklı F1 ile 82.8 HardCons değerine ulaşır. Hata analizi, ince ayarlı kodlayıcılarda yön tutarsızlığının küçük bir hata kaynağı olduğunu gösterir; asıl kaynak, tersinir ilişkileri tersinir olmayanlardan ayırt etmektir. Ayrıca SoftCons metriğinin sabit bir tahminciyle en yüksek değerine ulaştığını ve tersinir olmayan örneklerin yalnızca verinin çift yapısından tanınabildiğini gösteriyoruz.
\end{abstract}
